Floating photovoltaic (FPV) systems can expand solar generation without occupying the same land footprint as ground-mounted arrays. They can also share water-management and hydropower infrastructure where local conditions permit~\citep{sahu2016,lee2020}. Their practical value, however, depends on more than the solar resource over water. Ice, protected status, shoreline access, variable water area, grid connection, anchoring conditions, competing uses and permitting determine whether a resource estimate is useful for planning.

Existing assessments have established the scale and diversity of FPV opportunity. Spencer et al.~\citep{spencer2019} mapped technical potential on artificial water bodies in the United States. Lee et al.~\citep{lee2020} assessed global hydropower--FPV hybrids, whereas Jin et al.~\citep{jin2023} quantified reservoir electricity and water-saving scenarios. Woolway et al.~\citep{woolway2024} extended global analysis to lakes and linked generation with national electricity demand. National and regional studies have added protected-area, remote-sensing and economic screens at finer resolution~\citep{lopes2022,ates2022,faruqui2023,munozceron2023,jung2024,rosenlieb2025}. These totals answer different questions because their water-body classes, minimum sizes, coverage fractions and exclusion rules differ. Observed FPV coverage itself ranges from 0.004\% to 89.9\% across 643 mapped plants, with a mean of 34.2\%, which reinforces the need to present coverage as a scenario rather than a universal design rule~\citep{nobre2024}.

Three unresolved problems limit a progressively screened global estimate. First, global lake and reservoir inventories overlap but do not share complete identifiers. Proximity-only merging can either double count one water body or delete distinct neighbours~\citep{lehner2011,messager2016,wang2022geodar}. Second, evidence is uneven across water bodies and time. The public monthly series underlying the Global Lake Evaporation Volume (GLEV) dataset ends in 2018~\citep{zhao2022glev}; it therefore cannot reproduce a 1991--2020 dry-up criterion without additional data. Third, globally consistent infrastructure and engineering layers remain screening proxies. A road-density cell does not locate a shoreline route~\citep{meijer2018grip}; a mapped power line does not establish voltage, substation access or residual capacity; and modelled maximum water-body depth is not the depth at an array or anchor~\citep{khazaei2022globathy}. Missing evidence must consequently remain visible instead of becoming an implicit suitability pass.

We address these problems with an evidence-state assessment for HydroLAKES natural and controlled lakes of at least 1~km$^2$ and mapped GeoDAR/GRanD reservoirs of at least 0.01~km$^2$. We adjudicate lake--reservoir identity using authoritative and calibrated links, restore records removed without duplicate evidence, and hold the electricity model fixed while constraints are added. The analysis reports a reference physical resource (L0), climate-screened resource (L1), core and conservative public-evidence intervals (L2), and bounded road-and-power-line proximity intervals (partial L3). Coverage, ice source, surface-area criterion, infrastructure mapping, depth and cost are evaluated as separate sensitivities.

This structure supports two distinct planning decisions. Lower interval endpoints identify a globally comparable set with observed support for every applied public-data gate. Upper endpoints identify where targeted identity, surface-area or infrastructure data could most reduce uncertainty. Complete suitability under the published 1991--2020 criterion and complete project deployability are retained as unavailable states. The resulting sequence provides an applied screening framework without assigning project-level meaning to global proxies.
